\subsection{Estrutura de controle}
\label{subsec:malha}

\begin{figure}[htbp]
    \centering
    \begin{tikzpicture}[
        >=Latex,
        font=\small,
        block/.style={draw, thick, minimum height=9mm, align=center, fill=white},
        sum/.style={draw, circle, thick, inner sep=0pt, minimum size=7mm},
        wire/.style={->, thick},
        tap/.style={circle,fill=black,inner sep=1.25pt}
    ]
    % Sinais de entrada, erros e integrador discreto
    \node (ref) at (0,0) {$r(k)$};
    \node[sum] (s1) at (1.55,0) {};
    \node[sum] (s2) at (3.20,0) {};
    \node[block,minimum width=1.15cm] (delay) at (4.5,0) {$z^{-1}$};
    \node[block,minimum width=1.00cm] (ki) at (6.4,0) {$K_i$};
    \node[sum] (s3) at (7.75,0) {};
    \node[block,minimum width=3.6cm,minimum height=1.28cm] (plant) at (11.25,0) {\shortstack{Planta farmacocinética\\$x^{+}=Ax+Bu$\\$y=Cx$}};
    \node (out) at (14.8,0) {$y(k)$};
    
    \draw[wire] (ref) -- (s1);
    \draw[wire] (s1) -- node[above] {$e(k)$} (s2);
    \draw[wire] (s2) -- (delay);
    \draw[wire] (delay) -- (ki);
    \draw[wire] (ki) -- (s3);
    \draw[wire] (s3) -- node[above] {$u(k)$} (plant);
    \draw[wire] (plant) -- (out);
    
    % Sinais de soma
    \node at ($(s1)+(-0.15,0.0)$) {\scriptsize$+$};
    \node at ($(s1)+(0,-0.19)$) {\scriptsize$-$};
    \node at ($(s2)+(0.0,0.18)$) {\scriptsize$+$};
    \node at ($(s2)+(-0.15,0.0)$) {\scriptsize$+$};
    \node at ($(s3)+(-0.15,0.0)$) {\scriptsize$+$};
    \node at ($(s3)+(0,0.18)$) {\scriptsize$+$};
    
    % Realimentação do acumulador: saída de z^-1 volta ao somador 2
    \coordinate (accTap) at ($(delay.east)+(0.28,0)$);
    \node[tap] at (accTap) {};
    \draw[wire] (accTap) |- (3.2,1.25) -- (s2.north);
    
    % Realimentação proporcional pela saída medida
    \coordinate (yTap) at (13.75, 0);
    \node[tap] at (yTap) {};
    \node[block,minimum width=1.1cm] (kp) at (10.8,1.65) {$K_p$};
    \draw[wire] (yTap) |- (kp.east);
    \draw[wire] (kp.west) -- (7.75,1.65) -- (s3.north);
    
    % Realimentação negativa para o erro
    \draw[wire] (yTap) |- (1.55,-1.35) -- (s1.south);
    \end{tikzpicture}
    \caption{Malha de controle PI discreta para o modelo farmacocinético, sem os blocos de exportação de dados do Simulink.}
    \label{fig:malha_controle}
\end{figure}

De maneira geral, consideremos que $x(k)\in\mathbb{R}^{n_x}$, $u(k)\in\mathbb{R}^{n_u}$ e $y(k)\in\mathbb{R}^{n_y}$. Para as matrizes na representação de estados \eqref{eq:representacao_em_estados_v1} consideremos que $A(\alpha)\in\mathbb{R}^{n_x\times n_x}$, $B\in\mathbb{R}^{n_x\times n_u}$ e $C\in\mathbb{R}^{n_y\times n_x}$. Com o objetivo de rastrear uma referência de concentração $r(k)\in\mathbb{R}^{n_y}$, foi adotada uma malha fechada com ação proporcional e integral (Figura~\ref{fig:malha_controle}) cuja lei de controle é descrita por
\begin{equation}
 u(k)=K_p y(k)+K_i x_i(k).
 \label{eq:lei_controle}
\end{equation}
$K_p$ e $K_i$ são os ganhos a serem determinados pelo procedimento de projeto. O estado integral $x_i(k)$ é usado para eliminar o erro em regime permanente, e é definido por
\begin{equation}
 x_i(k+1)=x_i(k)+e(k)=x_i(k)+r(k)-y(k).
 \label{eq:integrador}
\end{equation}

\noindent
As duas variáveis de estado no tempo futuro podem ser obtidas substituindo a equação \eqref{eq:lei_controle} em \eqref{eq:planta_estados}, e substituindo a equação \eqref{eq:integrador}  em \eqref{eq:planta_saida}
\begin{subequations}\label{eq:planta_fechada_completa}
    \begin{align}
        x(k+1)=\bigl(A(\alpha)+BK_pC\bigr)x(k)+BK_i x_i(k),
        \label{eq:planta_fechada}\\
        x_i(k+1)=-Cx(k)+x_i(k)+r(k).
        \label{eq:integrador_fechado}
    \end{align}
\end{subequations}

Seja o vetor de estados aumentado $\widetilde{x}(k)=[x(k)^\top\;x_i(k)]^\top\in\mathbb{R}^{n_a}$ com $n_a=n_x+n_y$. Sob este novo vetor de estados, a dinâmica do sistema em malha fechada pode ser representada de maneira matricial, por
\begin{equation}
 \widetilde{x}(k+1)=
 \underbrace{\begin{bmatrix}
 A(\alpha)+BK_pC & BK_i\\
 -C & I_{n_y}
 \end{bmatrix}}_{A_r(\alpha)}\widetilde{x}(k)
 +\begin{bmatrix}0_{n_x\times n_y}\\I_{n_y}\end{bmatrix}r(k).
 \label{eq:sistema_aumentado}
\end{equation}
\noindent
Deste processo, obtem-se a matriz dinâmica do sistema aumentado $A_r(\alpha)\in\mathbb{R}^{n_a\times n_a}$.
